\section*{Capítulo \ref{ch:b3:virology} --- Virología}

\begin{solution}{exo:b3:virology:1}
Polio: IV, ARN de hebra positiva --- una ARN-polimerasa dependiente de
ARN. Gripe: V, ARN de hebra negativa --- su propia ARN-polimerasa,
llevada en el \omterm{def:b3:virology:virus}{virión}. VIH: VI --- transcriptasa inversa (e integrasa).
Herpes simple: I, ADN de doble hebra --- el hospedador podría en
principio aportarlo todo, pero trae su propia ADN-polimerasa y su
timidina-quinasa. Hepatitis B: VII --- una transcriptasa inversa para
copiar su ARN pregenómico en ADN. Rotavirus: III, ARN de doble hebra ---
una ARN-polimerasa dependiente de ARN dentro de la partícula.
\end{solution}

\begin{solution}{exo:b3:virology:2}
El ADN del fago ($^{32}$P) entró en las bacterias y la proteína
($^{35}$S) se quedó fuera y pudo arrancarse; la descendencia del fago
heredó el $^{32}$P. Así que el ADN es lo que el fago inyecta y lo que
dirige la fabricación de fago nuevo. Si la proteína hubiera llevado las
instrucciones, el $^{35}$S se habría encontrado dentro y habría pasado a
la descendencia.
\end{solution}

\begin{solution}{exo:b3:virology:3}
$84/(0.1\times 10^{-7}) = 8.4\times 10^{9}$ unidades formadoras de placa
por mL. El microscopio cuenta todas las partículas, incluidas las
\omterm{def:b3:virology:virus}{cápsides} vacías, las de genoma defectuoso y las que no consiguen unirse o
inyectar; solo las que completan una infección forman placas.
\end{solution}

\begin{solution}{exo:b3:virology:4}
Adhesión: inhibidores de entrada (maraviroc), anticuerpos
neutralizantes. Entrada y fusión: inhibidores de la fusión, bloqueantes
de la acidificación endosómica. Desensamblaje: fármacos que unen la
\omterm{def:b3:virology:virus}{cápside} (pleconaril; lenacapavir). Expresión y replicación: \omterm{def:b3:virology:drugs}{análogos de
nucleósido}, inhibidores de la transcriptasa inversa y de la integrasa,
interferón, \omterm{def:b3:rna-regulation:rnai}{interferencia por ARN}. Ensamblaje: \omterm{def:b3:virology:drugs}{inhibidores de proteasa}
(las poliproteínas sin cortar no pueden ensamblarse). Liberación:
inhibidores de la neuraminidasa (gripe) y linfocitos~T citotóxicos que
matan a la célula antes de la liberación.
\end{solution}

\begin{solution}{exo:b3:virology:5}
Una proteína de \qty{30}{kDa} son $30\,000/110 = 273$ residuos, $818$
nucleótidos: \qty{0.82}{kb}, el \qty{18}{\%} de un genoma de
\qty{4.5}{kb}. Codificar por separado las $180$ copias costaría
\qty{147}{kb}, treinta y tres veces el genoma entero; la propia proteína
de la \omterm{def:b3:virology:virus}{cápside} (\qty{5.4}{MDa}) pesa cuatro veces más que el genoma
(\qty{1.5}{MDa}).
\end{solution}

\begin{solution}{exo:b3:virology:6}
$V = 10^{5}(3e^{-0.5t} - 0.5e^{-3t})/2.5$: $t = 0.5$: $8.9\times
10^{4}$; $t = 2$: $4.4\times 10^{4}$; $t = 10$: $8\times 10^{2}$. Baja de
$50$ cuando $1.2\times 10^{5}e^{-0.5t} = 50$: $t = 2\ln(2400) \approx
\qty{16}{d}$.
\end{solution}

\begin{solution}{exo:b3:virology:7}
$uL = 0.36$; fracción sin error, $e^{-0.36} = 0.70$. Al triplicarse:
$uL =
1.08$, $e^{-1.08} = 0.34$ --- la mayor parte de la descendencia lleva ya
mutaciones y la secuencia más apta se diluye hacia el umbral. Un
coronavirus de \qty{30}{kb} con $u = 10^{-4}$ tendría $uL = 3$ y solo un
\qty{5}{\%} de copias sin error, más allá del umbral; su exonucleasa
lleva $u$ a $10^{-5}$, $uL = 0.3$, y tres cuartas partes sin error.
\end{solution}

\begin{solution}{exo:b3:virology:8}
Deriva: las mutaciones puntuales de la hemaglutinina y de la
neuraminidasa se acumulan bajo la selección de los anticuerpos, de modo
que la cepa de cada temporada es en parte nueva y la inmunidad del año
pasado está en parte obsoleta. Salto: la coinfección de una misma célula
por una cepa humana y otra animal reasocia los ocho segmentos, y un
subtipo de hemaglutinina de aves o de cerdos aparece en un \omterm{def:b3:virology:virus}{virus} adaptado
al ser humano. Nadie tiene anticuerpos contra el subtipo nuevo, de modo
que toda la población es susceptible a la vez --- la condición para una
pandemia, cumplida en 1918, 1957, 1968 y 2009, cada vez con una
hemaglutinina nueva.
\end{solution}

\begin{solution}{exo:b3:virology:9}
La timidina-quinasa del herpesvirus fosforila el aciclovir, al que las
quinasas celulares apenas tocan; las quinasas celulares completan después
el trifosfato, que la ADN-polimerasa vírica incorpora y que, al carecer
de hidroxilo 3$'$, termina la cadena. Las células no infectadas nunca
fabrican la forma activa, y la polimerasa vírica la prefiere a la
celular. Resistencia: la pérdida o la alteración de la timidina-quinasa
vírica (lo más frecuente), o mutaciones de la polimerasa vírica que
excluyan el análogo.
\end{solution}

\begin{solution}{exo:b3:virology:10}
Con $T$ constante las ecuaciones son lineales en $(I, V)$ con matriz
$\begin{pmatrix} -\delta & \beta T \\ p & -c \end{pmatrix}$, traza
$-(\delta + c) < 0$ y determinante $\delta c - \beta Tp$. Los valores
propios tienen traza negativa, de modo que uno es positivo exactamente
cuando el determinante es negativo: $\beta Tp > c\delta$, es decir,
$R_{0} = \beta
Tp/(c\delta) > 1$. Factores: $p/\delta$ es el número de viriones que
fabrica una célula infectada a lo largo de su vida; $\beta T/c$ es el
número de células que infecta un \omterm{def:b3:virology:virus}{virión} en la suya; y su producto es el
número de células infectadas que produce una célula infectada.
\end{solution}

\begin{solution}{exo:b3:virology:11}
$10^{6} = 2^{20}$: veinte semividas, $880$ meses, unos $73$ años --- más
que una vida. Los fármacos que bloquean la replicación no tocan un
provirus que no se está transcribiendo; el reservorio persiste y, cuando
se interrumpe el tratamiento, basta una sola célula que se reactive para
reiniciar la infección. Una cura ha de eliminar o silenciar de forma
permanente el reservorio: matar las células latentes (reactivarlas bajo
tratamiento para que mueran o las maten), escindir o inutilizar el
provirus, o sustituir el sistema inmunitario por uno al que el \omterm{def:b3:virology:virus}{virus} no
pueda infectar.
\end{solution}

\begin{solution}{exo:b3:virology:12}
$uL = 0.03$ mutaciones por genoma; $10^{9}$ genomas al día llevan
$3\times
10^{7}$ mutaciones entre $9\times 10^{4}$ mutaciones puntuales posibles:
cada una surge unas $300$ veces al día. Un doble mutante concreto surge a
$10^{-12}$ por genoma, $10^{-3}$ al día --- una vez cada tres años; y uno
triple, a $10^{-18}$, nunca. Dos fármacos con mutaciones de resistencia
independientes bastarían en principio, y tres dan margen para una mala
adherencia. Dos fármacos que se unan a la misma enzima pueden ser
derrotados por una sola mutación que altere el sitio o la conformación de
la enzima, y sus mutaciones de resistencia no son independientes; cuentan
como menos de dos.
\end{solution}

\begin{solution}{pb:b3:virology:1}
\textbf{1.} $2\times 10^{5}\times 1.5\times 10^{4} = 3\times 10^{9}$
viriones.
\textbf{2.} Aclarados, $cV = 3\times 3\times 10^{9} \approx 10^{10}$ al
día; producidos, otros tantos.
\textbf{3.} $I^{*} = cV^{*}/p = 9\times 10^{9}/500 = 1.8\times 10^{7}$
células infectadas; mueren $\delta I^{*} = 9\times 10^{6}$ al día.
\textbf{4.} ARN, $2\times 9700\times 330 = 6.4\times 10^{6}$ Da;
proteína, $2000\times 25\,000 = 5\times 10^{7}$ Da; total,
$5.6\times 10^{7}$ Da $\approx 10^{-16}$ g. Todos los viriones juntos:
$3\times 10^{9}\times
10^{-16} = 3\times 10^{-7}$ g, un tercio de microgramo --- tres mil veces
menos que un grano de arena.
\textbf{5.} $1.8\times 10^{7}/10^{12} = 2\times 10^{-5}$ de la reserva.
Diez años de $9\times 10^{6}$ muertes al día son $3\times 10^{10}$
células, el tres por ciento de la reserva: la reserva falla no por una
resta aritmética, sino porque una década de regeneración forzada y de
activación crónica agota la capacidad de reponerla.
\textbf{6.} $9\times 10^{6}/24 \approx 4\times 10^{5}$ cadáveres
eliminándose en cada momento.
\textbf{7.} $V = 2\times 10^{5}(3e^{-0.5t} - 0.5e^{-3t})/2.5$: $t = 1$:
$1.4\times 10^{5}$; $t = 3$: $5.4\times 10^{4}$; $t = 7$: $7\times
10^{3}$; $t = 14$: $2\times 10^{2}$.
\textbf{8.} $2.4\times 10^{5}e^{-0.5t} = 50$: $t = 2\ln(4800) \approx 17$
días.
\textbf{9.} Los días 3--10 caen en la fase lenta: la pendiente de $\ln V$
ahí es $-\delta$ (de $5.4\times 10^{4}$ a $1.6\times 10^{3}$ en siete
días, $\delta = 0.50$). La fase rápida dura solo horas ($1/c$), de modo
que $c$ exige muestras cada pocas horas del primer día, ajustadas a la
fórmula de dos exponenciales.
\textbf{10.} Las células infectadas de forma latente, que liberan \omterm{def:b3:virology:virus}{virus}
al reactivarse; su velocidad de decaimiento es $\ln 2/44$ por mes, unos
$5\times 10^{-4}$ por día.
\textbf{11.} \omterm{def:b3:virology:virus}{Virión}, $\ln 2/3 = \qty{0.23}{d}$, unas \qty{5.5}{h};
célula infectada, $\ln 2/0.5 = \qty{1.4}{d}$.
\textbf{12.} Una carga constante parecía un \omterm{def:b3:virology:virus}{virus} que no hacía nada; las
ecuaciones la muestran como un estado estacionario en el que se fabrican
y se aclaran $10^{10}$ viriones y se infectan y mueren $10^{7}$
linfocitos~T cada día.
\textbf{13.} $9700\times 10^{-5} \approx 0.1$ mutaciones por genoma;
$10^{10}\times 0.1 = 10^{9}$ mutaciones al día.
\textbf{14.} $3\times 9700 \approx \num{29000}$ mutaciones puntuales
posibles; cada una surge $10^{9}/\num{29000} \approx 3\times 10^{4}$
veces al día.
\textbf{15.} Doble concreto: $10^{-10}$ por genoma, alrededor de uno al
día sin tratamiento; triple concreto: $10^{-15}$, $10^{-5}$ al día ---
una vez cada tres siglos.
\textbf{16.} $10^{7}\times 10^{-15} = 10^{-8}$ al día, $4\times 10^{-6}$
al año.
\textbf{17.} Con un fármaco fuera, bastan los mutantes resistentes a los
otros dos: un doble concreto a $10^{-10}$; $7\times 10^{7}$ genomas en la
semana dan $7\times 10^{-3}$ esperados --- todavía improbable, salvo que
la carga rebote hacia $10^{10}$ al día durante el lapso, cuando pasa a
ser uno al día y el escape es probable. Las dosis olvidadas son la vía
por la que llega la resistencia.
\textbf{18.} Los análogos comparten el sitio activo de la transcriptasa
inversa, y una sola mutación allí (o un conjunto de mutaciones frente a
análogos de timidina) altera el manejo que hace la enzima de varios de
ellos: la resistencia cruzada. Dos nucleósidos cuentan como menos de dos
fármacos independientes, de modo que las pautas combinan clases --- un
nucleósido, un inhibidor de la integrasa, un inhibidor de la proteasa o
un no nucleósido.
\textbf{19.} $20$ semividas, $880$ meses, $73$ años; hasta $10^{4}$,
$\log_{2}100 = 6.6$ semividas, $290$ meses, $24$ años.
\textbf{20.} De un \omterm{def:b3:virology:virus}{virión} a $3\times 10^{9}$ con $R_{0} = 8$ por
generación: $\ln(3\times 10^{9})/\ln 8 = 10.5$ generaciones, tres
semanas.
\textbf{21.} Sacudir y matar: reactivar las células latentes bajo
tratamiento para que mueran o las maten --- obstáculo: ningún agente las
reactiva todas, y las reactivadas no mueren de forma fiable. Edición
génica: escindir el provirus o destruir el receptor CCR5 en las células
del paciente --- obstáculo: llegar a todas las células. (Sustituir la
médula por células de un donante sin CCR5 ha curado a un puñado de
pacientes, con un riesgo aceptable solo cuando hacía falta un trasplante
de todos modos.)
\textbf{22.} $1 - 1/R_{0}$: el \qty{67}{\%} para $R_{0} = 3$ y el
\qty{80}{\%} para $5$.
\textbf{23.} El tratamiento lleva $\beta \to 0$ dentro del paciente, la
carga cae por el teorema hasta casi nada y la transmisión, que es
proporcional a la carga, se detiene; tratar a todas las personas
infectadas lleva el $R_{0}$ de la población por debajo de uno.
\textbf{24.} $f = (1 - 1/R_{0})/E$: con $R_{0} = 4$ y $E = 0.7$, $f =
0.75/0.7 = 1.07$ --- más que todo el mundo: la \omterm{def:b3:virology:drugs}{vacuna} por sí sola no
puede alcanzar el umbral; con $E = 0.9$, $f = 0.83$.
\textbf{25.} Unos $10^{10}$ viriones producidos al día; indetectable
hacia el día $17$; y unos $4\times 10^{-6}$ triples mutantes al año con
tratamiento.
\end{solution}
